% Volume VIII: Computational Neuroscience Mathematics
% Continuity note for Chapter 46.
% Previous dependencies: Chapters 13-15 define random variables, distributions, likelihoods, and Bayes; Chapter 27 covers stochastic processes, Poisson processes, SDE notation, and filtering; Chapters 43-45 define voltage, spikes, synaptic filters, and plasticity.
% New durable concepts: spike trains as point processes, counting processes, conditional intensity, homogeneous and inhomogeneous Poisson processes, interarrival times, event-time likelihoods, neural GLMs, stimulus/history/coupling filters, maximum-likelihood estimation, regularization, time-rescaling, KS diagnostics, and held-out model checking.
% Forward hooks: Chapter 47 uses spike-train models for neural coding and information; Chapter 48 uses point-process observation models for latent dynamics; Chapter 51 uses likelihood and model comparison for neural data analysis.
% Notation commitments: spike times are 0<t_1<...<t_n<=T; counting processes are N(t); increments are dN(t) or N(t+dt)-N(t); histories are H_t; conditional intensity is lambda(t | H_t); integrated intensity is Lambda(a,b).
\section{Chapter 46: Spike Trains, Point Processes, and Neural GLMs}
\label{ch:spike-trains-point-processes-neural-glms}

Biophysical and point-neuron models generate spikes. Experimental neuroscience often starts from the other direction: a recording gives spike times, and we need a statistical model for those events. The events are discrete, the time axis is continuous or finely binned, and the probability of a spike usually depends on stimulus, recent spiking history, and activity in other neurons.

\paragraph{Chapter problem.}
How can spike trains be represented as point processes, how do Poisson and GLM models assign likelihoods to event times, and how can we check whether a fitted spike model is adequate?

\paragraph{Section map.}
Section 46.1 defines spike trains, counting processes, conditional intensity, homogeneous Poisson processes, and interarrival times. Section 46.2 extends to inhomogeneous Poisson models and event-time likelihoods. Section 46.3 builds neural generalized linear models with stimulus, history, and coupling filters. Section 46.4 develops model checking through held-out likelihood, residuals, and the time-rescaling theorem.

\subsection{46.1 Spike trains as point processes}

\paragraph{Guiding question.}
What is the right probability object for a sequence of spike times?

\paragraph{Beginner-facing intuition.}
A spike train is not a continuous voltage trace and not just a count. It is a set of event times. A point process is a stochastic process whose sample path is a collection of points on a time axis. The central object is not the probability that \(V(t)\) has some value, but the probability that an event occurs in a small time interval.

The conditional intensity is the point-process analogue of an instantaneous firing rate. It says how likely a spike is right now, given everything the model knows about the past.

\paragraph{Formal definitions and notation.}
Fix an observation window \([0,T]\). A spike train is an ordered set of event times
\[
0<t_1<t_2<\cdots<t_n\leq T.
\]
The associated counting process is
\[
N(t)=\#\{k:t_k\leq t\},\qquad 0\leq t\leq T.
\]
Its increment over a small interval is
\[
dN(t)=N(t+dt)-N(t),
\]
which is usually \(0\) or \(1\) for very small \(dt\).

Let \(\mathcal H_t\) denote the history available just before time \(t\): previous spikes, stimulus values up to \(t\), and any covariates included in the model. The conditional intensity is
\[
\boxed{
\lambda(t\mid\mathcal H_t)
=\lim_{dt\downarrow0}
\frac{P(dN(t)=1\mid\mathcal H_t)}{dt}
}
\]
when the limit exists. Equivalently,
\[
P(dN(t)=1\mid\mathcal H_t)=\lambda(t\mid\mathcal H_t)\,dt+o(dt),
\]
and
\[
P(dN(t)\geq2\mid\mathcal H_t)=o(dt).
\]
The units of \(\lambda\) are spikes per unit time.

\paragraph{Core derivation: exponential interarrival times for homogeneous Poisson spiking.}
A homogeneous Poisson process has constant intensity
\[
\lambda(t\mid\mathcal H_t)=r.
\]
Let \(T_1\) be the waiting time to the first spike. Divide time into small bins of width \(dt\). The probability of no spike in one bin is approximately \(1-r\,dt\). The probability of no spike before time \(t\) is approximately
\[
(1-r\,dt)^{t/dt}.
\]
Taking \(dt\to0\) gives
\[
P(T_1>t)=e^{-rt}.
\]
Therefore the density of \(T_1\) is
\[
f_{T_1}(t)=r e^{-rt},\qquad t\geq0.
\]
The interarrival times of a homogeneous Poisson process are independent exponential random variables with mean \(1/r\). This connects event-time and count views of the same process.

\paragraph{Concrete example.}
If a neuron is modeled as a homogeneous Poisson process with \(r=20\) Hz, then the expected waiting time between spikes is
\[
1/r=0.05\ \mathrm{s}=50\ \mathrm{ms}.
\]
In a \(100\) ms window, the expected count is \(rT=2\). The probability of zero spikes is
\[
P(N(0.1)=0)=e^{-2}\approx0.135.
\]
The probability of exactly two spikes is
\[
e^{-2}\frac{2^2}{2!}\approx0.271.
\]

\paragraph{Computational neuroscience example.}
Raster plots show repeated spike-train sample paths across trials. The counting process \(N(t)\) turns each raster row into a nondecreasing step function. The intensity \(\lambda(t\mid\mathcal H_t)\) is a model-based object: it may depend on a stimulus, a refractory effect, a network state, or a latent variable. A neuron with a refractory period cannot be exactly homogeneous Poisson because its conditional intensity is reduced immediately after a spike.

\paragraph{AI and machine-learning example.}
Temporal point processes are used for event streams such as clicks, transactions, messages, failures, and medical events. Neural temporal point-process models, Hawkes processes, and transformer-based event models all generalize the same idea: predict an intensity for the next event given history. Spike trains are one biologically important instance of this broader event-modeling problem.

\paragraph{Common misconceptions and failure modes.}
First, a spike train is not fully described by its mean rate; timing structure and history dependence can matter. Second, a point-process intensity is conditional on a model's chosen history, not a directly observed physical force. Third, Poisson does not mean ``random-looking''; it means a specific independent-increment process with exponential waiting times in the homogeneous case. Fourth, binning spike trains too coarsely can hide event timing and create bins with multiple spikes, violating small-bin approximations.

\paragraph{Exercises.}
\begin{enumerate}[label=\textbf{46.1.\arabic*.}, leftmargin=*]
\item \textbf{Mechanical.} Given spike times \(0.12,0.31,0.50\) seconds in a \(1\) second trial, write \(N(0.2)\), \(N(0.4)\), and \(N(1.0)\).
\item \textbf{Proof or derivation.} Derive the exponential waiting-time survival \(P(T_1>t)=e^{-rt}\) from the small-bin approximation.
\item \textbf{Numerical/computational.} For a homogeneous Poisson process at \(5\) Hz, compute the mean interspike interval and the probability of no spikes in \(400\) ms.
\item \textbf{Interpretation.} Why does a refractory period require history dependence in the intensity?
\item \textbf{Misconception/debugging.} A student says, ``The neuron fired two spikes in one trial, so its firing rate is exactly \(2/T\).'' Explain the difference between an observed count rate and the model intensity.
\end{enumerate}

\subsection{46.2 Inhomogeneous Poisson models}

\paragraph{Guiding question.}
How do we model spikes whose rate changes over time because of a stimulus or task variable?

\paragraph{Beginner-facing intuition.}
A homogeneous Poisson process has one constant rate. Real neural responses often have stimulus-locked time courses: a visual neuron may fire strongly after stimulus onset, adapt, and then return to baseline. An inhomogeneous Poisson process keeps Poisson randomness but allows the intensity \(\lambda(t)\) to be a deterministic or covariate-determined function of time.

The key quantity is integrated intensity. The expected number of events in a window is the area under the intensity curve over that window.

\paragraph{Formal definitions and notation.}
An inhomogeneous Poisson process on \([0,T]\) has deterministic nonnegative intensity
\[
\lambda\colon [0,T]\to\mathbb{R}_{\geq 0}.
\]
For any interval \([a,b]\subseteq[0,T]\), define the integrated intensity
\[
\boxed{\Lambda(a,b)=\int_a^b\lambda(t)\,dt.}
\]
The count in that interval satisfies
\[
N(b)-N(a)\sim\operatorname{Poisson}(\Lambda(a,b)).
\]
Counts in disjoint intervals are independent. Thus
\[
P(N(b)-N(a)=k)=e^{-\Lambda(a,b)}\frac{\Lambda(a,b)^k}{k!}.
\]

Given observed event times \(t_1,\ldots,t_n\) in \([0,T]\), the event-time likelihood is
\[
\boxed{
L(\lambda;t_1,\ldots,t_n)
=\left[\prod_{i=1}^n\lambda(t_i)\right]
\exp\!\left(-\int_0^T\lambda(t)\,dt\right).
}
\]
The log-likelihood is
\[
\ell(\lambda)=\sum_{i=1}^n\log\lambda(t_i)-\int_0^T\lambda(t)\,dt.
\]

\paragraph{Core derivation: likelihood from small bins.}
Partition \([0,T]\) into small bins of width \(\Delta\). In bin \(j\), approximate the spike probability by \(\lambda(t_j)\Delta\). If a bin contains one spike, its contribution is approximately \(\lambda(t_j)\Delta\). If it contains no spike, its contribution is approximately \(1-\lambda(t_j)\Delta\). Multiplying over bins gives
\[
\prod_{\text{spike bins}}\lambda(t_j)\Delta
\prod_{\text{empty bins}}(1-\lambda(t_j)\Delta).
\]
As \(\Delta\to0\), the product over empty bins approaches
\[
\exp\!\left(-\int_0^T\lambda(t)\,dt\right).
\]
The product over spike bins gives \(\prod_i\lambda(t_i)\) times a factor depending on bin width and event ordering. For likelihood comparisons over \(\lambda\), the relevant density is the boxed expression above.

\paragraph{Concrete example.}
Consider a one-second trial with piecewise constant intensity:
\[
\lambda(t)=
\begin{cases}
5\ \mathrm{Hz}, & 0\leq t<0.5,\\
25\ \mathrm{Hz}, & 0.5\leq t\leq1.
\end{cases}
\]
The integrated intensity is
\[
\Lambda(0,1)=5(0.5)+25(0.5)=15.
\]
The expected spike count is \(15\). The probability of exactly \(12\) spikes is
\[
e^{-15}\frac{15^{12}}{12!}\approx0.082.
\]
If two models have the same total \(\Lambda\) but assign different intensity at the observed spike times, the product \(\prod_i\lambda(t_i)\) distinguishes them.

\paragraph{Computational neuroscience example.}
The peri-stimulus time histogram (PSTH) estimates a time-varying firing rate by averaging binned spike counts across repeated trials aligned to stimulus onset. An inhomogeneous Poisson model treats the PSTH-like rate as the trial-shared intensity and assumes remaining trial-to-trial variability is Poisson. Deviations from this model can indicate refractoriness, bursting, gain fluctuations, or hidden state changes.

\paragraph{AI and machine-learning example.}
In event prediction, an inhomogeneous Poisson model can use a known covariate such as time of day, stimulus intensity, or user exposure to set \(\lambda(t)\). It is the point-process analogue of a nonstationary baseline model. More complex neural point-process models replace the hand-chosen \(\lambda(t)\) by a learned function of covariates and history.

\paragraph{Common misconceptions and failure modes.}
First, a time-varying intensity alone does not create spike-history dependence; disjoint intervals remain independent given \(\lambda(t)\). Second, the integrated intensity is not the same as the instantaneous intensity. It is an expected count. Third, fitting a smooth rate to one trial is statistically fragile; repeated trials or structural assumptions are needed. Fourth, a good PSTH fit does not guarantee correct spike timing statistics.

\paragraph{Exercises.}
\begin{enumerate}[label=\textbf{46.2.\arabic*.}, leftmargin=*]
\item \textbf{Mechanical.} Compute \(\Lambda(0,2)\) for \(\lambda(t)=3+2t\) Hz on a two-second interval.
\item \textbf{Proof or derivation.} Derive the event-time log-likelihood \(\sum_i\log\lambda(t_i)-\int_0^T\lambda(t)\,dt\) from small-bin Bernoulli probabilities.
\item \textbf{Numerical/computational.} For \(\Lambda=4\), compute the probability of \(0\), \(1\), and \(2\) spikes in the window.
\item \textbf{Interpretation.} A model predicts the right total spike count but places high intensity before the observed response peak. Which term in the likelihood penalizes this?
\item \textbf{Misconception/debugging.} A student says, ``Because \(\lambda(t)\) changes over time, the process has memory.'' Explain why time variation and history dependence are different.
\end{enumerate}

\subsection{46.3 Neural generalized linear models}

\paragraph{Guiding question.}
How can stimulus features, spike history, and other neurons' activity be combined into a fitted conditional intensity model?

\paragraph{Beginner-facing intuition.}
A neural GLM is a regression model for spikes. It builds a linear predictor from covariates, then passes that predictor through a positive nonlinearity to get an intensity. The covariates can include stimulus filters, post-spike history filters, and coupling filters from other neurons. Because the model has a likelihood, its parameters can be estimated and compared statistically.

\paragraph{Formal definitions and notation.}
Discretize time into bins of width \(\Delta\). Let \(y_t\in\{0,1,2,\ldots\}\) be the spike count in bin \(t\). Let \(x_t\in\mathbb{R}^p\) be a design vector containing stimulus features, recent spike-history values, coupling covariates, and a constant \(1\) for the bias. Let \(\theta\in\mathbb{R}^p\) be model parameters.

A Poisson GLM uses
\[
\eta_t=\theta^\top x_t,
\qquad
\lambda_t=\exp(\eta_t),
\]
and
\[
y_t\mid x_t,\mathcal H_t\sim\operatorname{Poisson}(\Delta\lambda_t).
\]
The binned log-likelihood is
\[
\boxed{
\ell(\theta)=
\sum_t\left[
y_t\log\Delta+y_t\,\theta^\top x_t-\Delta e^{\theta^\top x_t}-\log(y_t!)
\right].
}
\]

A common neural design is
\[
\eta_t=b+k^\top s_t+h^\top y_{t-1:t-L}
+\sum_{j\neq i} c_j^\top y^{(j)}_{t-1:t-L},
\]
where \(k\) is a stimulus filter, \(h\) is a self-history filter, and \(c_j\) are coupling filters from other neurons.

\paragraph{Core derivation: GLM score equation.}
Let \(X\in\mathbb{R}^{T\times p}\) be the design matrix with rows \(x_t^\top\), let \(y\in\mathbb{R}^T\), and let \(\lambda_t=e^{\theta^\top x_t}\). Ignoring constants \(y_t\log\Delta-\log(y_t!)\), which are independent of \(\theta\),
\[
\tilde\ell(\theta)=\sum_t y_t\theta^\top x_t-\Delta\sum_t e^{\theta^\top x_t}.
\]
Differentiate:
\[
\nabla_\theta\ell(\theta)=\nabla_\theta\tilde\ell(\theta)
=\sum_t y_t x_t-\Delta\sum_t e^{\theta^\top x_t}x_t.
\]
Thus
\[
\boxed{\nabla_\theta\ell(\theta)=X^\top(y-\Delta\lambda).}
\]
At a maximum-likelihood estimate without regularization, the fitted expected sufficient statistics \(X^\top\Delta\lambda\) match the observed statistics \(X^\top y\), subject to the model class.

\paragraph{Concrete example.}
Suppose \(\Delta=0.01\) s and a model has only a bias and one stimulus feature:
\[
\eta_t=b+k s_t.
\]
Let \(b=\log 20\), \(k=0.5\), and \(s_t=2\). Then
\[
\lambda_t=e^{\log20+1}=20e\approx54.4\ \mathrm{Hz}.
\]
The expected count in a \(10\) ms bin is
\[
\Delta\lambda_t=0.01(54.4)=0.544.
\]
The probability of at least one spike in the bin under the Poisson approximation is
\[
1-e^{-0.544}\approx0.420.
\]
If \(s_t=0\), the intensity is \(20\) Hz and the expected bin count is \(0.2\).

\paragraph{Computational neuroscience example.}
A retinal ganglion cell GLM may use a spatiotemporal stimulus filter \(k\), a spike-history filter \(h\), and coupling filters from nearby cells. A negative self-history filter immediately after a spike captures refractoriness. A later positive lobe can capture bursting. Coupling filters can reveal statistical dependencies, but they do not by themselves prove direct synaptic connectivity.

\paragraph{AI and machine-learning example.}
The neural GLM is a transparent ancestor of modern sequence models for events. It has a linear feature map, an exponential link, and a log-likelihood trained by gradient methods. Deep point-process models replace \(\theta^\top x_t\) with a neural network hidden state, but the intensity, likelihood, and model-checking ideas remain point-process ideas.

\paragraph{Common misconceptions and failure modes.}
First, GLM filters are model parameters, not necessarily biological filters in a one-to-one sense. Second, coupling filters can reflect shared stimulus drive or common latent state, not only monosynaptic connections. Third, the exponential link guarantees positivity but can produce unrealistically large rates without regularization or sufficient data. Fourth, a binned Bernoulli GLM and a binned Poisson GLM are close only when bins are small enough that multiple spikes are rare.

\paragraph{Exercises.}
\begin{enumerate}[label=\textbf{46.3.\arabic*.}, leftmargin=*]
\item \textbf{Mechanical.} In \(\eta_t=b+k^\top s_t+h^\top y_{t-1:t-L}\), identify the baseline, stimulus, and history terms.
\item \textbf{Proof or derivation.} Derive \(\nabla_\theta\ell=X^\top(y-\Delta\lambda)\) for the Poisson GLM.
\item \textbf{Numerical/computational.} For \(\Delta=0.005\), \(b=\log 10\), \(k=1\), and \(s_t=-1\), compute \(\lambda_t\), the expected bin count, and the probability of zero spikes in the bin.
\item \textbf{Interpretation.} What does a strongly negative spike-history filter immediately after a spike usually represent?
\item \textbf{Misconception/debugging.} A fitted coupling filter from neuron A to neuron B is positive. Explain two reasons this need not prove a direct excitatory synapse.
\end{enumerate}

\subsection{46.4 Model checking for spike models}

\paragraph{Guiding question.}
After fitting a spike model, how can we tell whether it captures the event statistics that matter?

\paragraph{Beginner-facing intuition.}
Fitting maximizes a criterion within a model class. It does not prove the model is correct. A spike model may fit average rates but miss refractoriness, overdispersion, bursts, stimulus timing, or cross-neuron dependencies. Model checking asks whether the residual structure left behind looks like what the model assumes.

For point processes, the most important diagnostic idea is time rescaling. If the conditional intensity is correct, then time measured in units of predicted cumulative intensity should turn the spike train into a unit-rate Poisson process.

\paragraph{Formal definitions and notation.}
Let a fitted model provide conditional intensity \(\hat\lambda(t\mid\mathcal H_t)\). For consecutive spike times \(t_{i-1}<t_i\), define transformed intervals
\[
z_i=\int_{t_{i-1}}^{t_i}\hat\lambda(u\mid\mathcal H_u)\,du.
\]
If the model is correct under regularity assumptions, then
\[
z_1,z_2,\ldots
\]
are independent \(\operatorname{Exponential}(1)\) random variables. Equivalently,
\[
u_i=1-e^{-z_i}
\]
should be independent \(\operatorname{Uniform}(0,1)\) random variables.

Other diagnostics include held-out log-likelihood,
\[
\ell_{\mathrm{test}}=\sum_{i\in\mathrm{test}}\log\hat\lambda(t_i)-\int_{\mathrm{test}}\hat\lambda(t)\,dt,
\]
residual count checks, calibration by condition, Fano factor comparisons, and posterior predictive simulations.

\paragraph{Core theorem: time-rescaling intuition.}
For an inhomogeneous Poisson process with deterministic intensity \(\lambda(t)\), define transformed time
\[
\tau(t)=\int_0^t\lambda(u)\,du.
\]
In transformed time, the expected number of spikes between \(\tau=a\) and \(\tau=b\) is \(b-a\). The probability of no spike between original times \(s\) and \(t\) is
\[
\exp\!\left(-\int_s^t\lambda(u)\,du\right).
\]
Therefore the transformed waiting time
\[
z=\int_s^t\lambda(u)\,du
\]
has survival \(P(z>c)=e^{-c}\), the exponential survival with rate \(1\). Conditional-intensity models use the same idea after conditioning on history: if the conditional intensity is right, then accumulated conditional hazard between spikes behaves like exponential waiting time.

\paragraph{Concrete example.}
Suppose a fitted model predicts constant \(\hat\lambda=20\) Hz, and observed spike times are \(0.10\), \(0.16\), and \(0.31\) seconds. The transformed intervals are
\[
z_1=20(0.10-0)=2.0,
\]
\[
z_2=20(0.16-0.10)=1.2,
\]
\[
z_3=20(0.31-0.16)=3.0.
\]
The corresponding uniform variables are
\[
u_i=1-e^{-z_i},
\]
giving approximately \(0.865\), \(0.699\), and \(0.950\). With many intervals, these values can be compared with a uniform distribution using a KS plot. A systematic excess near \(1\) would suggest intervals are longer than predicted, meaning the model overpredicts spiking between observed events.

\paragraph{Computational neuroscience example.}
For a retinal GLM, one should check held-out likelihood, PSTH prediction, spike-history residuals, and time-rescaled intervals. A model with no history term may match stimulus-locked firing rate but fail time-rescaling because it predicts too many spikes immediately after an observed spike. Adding a refractory history filter can fix that specific failure.

\paragraph{AI and machine-learning example.}
Modern event models are often evaluated by held-out negative log-likelihood, calibration of event times, and simulation-based checks. Time-rescaling is the point-process analogue of probability calibration in classifiers: predicted probabilities should match empirical frequencies after conditioning on the model information. A high-capacity neural model can still be miscalibrated even if its training likelihood is high.

\paragraph{Common misconceptions and failure modes.}
First, high training likelihood is not model checking; it may reflect overfitting. Second, matching mean rate is not enough for a point-process model. Third, a KS plot checks a particular transformation and can miss structured dependencies unless independence is also examined. Fourth, failure of a diagnostic does not automatically identify the biological mechanism; it tells which statistical assumption is violated.

\paragraph{Exercises.}
\begin{enumerate}[label=\textbf{46.4.\arabic*.}, leftmargin=*]
\item \textbf{Mechanical.} Define the transformed interval \(z_i\) used in the time-rescaling diagnostic.
\item \textbf{Proof or derivation.} Show that if \(Z\sim\operatorname{Exponential}(1)\), then \(U=1-e^{-Z}\) is uniform on \([0,1]\).
\item \textbf{Numerical/computational.} A model predicts \(\hat\lambda(t)=10\) Hz. Observed intervals are \(50\) ms, \(100\) ms, and \(300\) ms. Compute the transformed intervals \(z_i\).
\item \textbf{Interpretation.} A GLM matches the PSTH but its rescaled intervals are too small immediately after spikes. What model component is likely missing?
\item \textbf{Misconception/debugging.} A student says, ``The test likelihood is better than baseline, so the model is true.'' Explain why model comparison and absolute goodness-of-fit are different.
\end{enumerate}

\paragraph{Closing bridge.}
Point-process models turn spike trains into likelihood-bearing mathematical objects. Chapter 47 uses those models to ask what neural responses encode: how tuning curves, population codes, decoders, and information measures connect stimuli to spikes and spikes back to inferred variables.

\clearpage
